\section*{Capitulo 1. Numeros reales y herramientas numericas}
\addcontentsline{toc}{section}{Capitulo 1. Numeros reales y herramientas numericas}

\begin{shortanswer}{[GX-C01.S01-01]}
\(-2\in\mathbb{Z}\).
\end{shortanswer}

\begin{shortanswer}{[GX-C01.S01-02]}
\(0.\overline{3}\in\mathbb{Q}\).
\end{shortanswer}

\begin{shortanswer}{C01.S01 practica autonoma}
\begin{enumerate}[label=\textbf{PX-C01.S01-\arabic*:}, leftmargin=*]
  \item \(\mathbb{R}\setminus\mathbb{Q}\)
  \item \(\mathbb{Q}\)
  \item \(\mathbb{Z}\)
  \item \(\mathbb{Z}\)
  \item \(\mathbb{Z}\)
  \item \(\mathbb{Q}\)
\end{enumerate}
\end{shortanswer}

\begin{shortanswer}{[GX-C01.S02-01]}
\(\dfrac{2}{11}\).
\end{shortanswer}

\begin{shortanswer}{[GX-C01.S02-02]}
\(\dfrac{49}{20}\).
\end{shortanswer}

\begin{shortanswer}{C01.S02 practica autonoma}
\begin{enumerate}[label=\textbf{PX-C01.S02-\arabic*:}, leftmargin=*]
  \item \(\dfrac{13}{18}\)
  \item \(\dfrac{32}{9}\)
  \item \(\dfrac{1}{8}\)
  \item \(\dfrac{31}{15}\)
  \item Decimal exacto
  \item Decimal periodico mixto
\end{enumerate}
\end{shortanswer}

\begin{shortanswer}{[GX-C01.S03-01]}
\((-\infty,-1)\cup [4,+\infty)\).
\end{shortanswer}

\begin{shortanswer}{[GX-C01.S03-02]}
\((1,5)\).
\end{shortanswer}

\begin{shortanswer}{C01.S03 practica autonoma}
\begin{enumerate}[label=\textbf{PX-C01.S03-\arabic*:}, leftmargin=*]
  \item \((-2,3]\)
  \item \(\{1\}\)
  \item \([0,8)\)
  \item \((2,+\infty)\)
  \item \(\varnothing\)
  \item \((-\infty,3)\)
\end{enumerate}
\end{shortanswer}

\begin{shortanswer}{[GX-C01.S04-01]}
\(x=2\) o \(x=-4\).
\end{shortanswer}

\begin{shortanswer}{[GX-C01.S04-02]}
\(-2<x<6\).
\end{shortanswer}

\begin{shortanswer}{C01.S04 practica autonoma}
\begin{enumerate}[label=\textbf{PX-C01.S04-\arabic*:}, leftmargin=*]
  \item \(x=\pm 7\)
  \item \(x=1\) o \(x=-\dfrac{1}{3}\)
  \item \([-5,-3]\)
  \item \(x=3\) o \(x=7\)
  \item \(x<-2\) o \(x>-1\)
  \item \(x=1\)
\end{enumerate}
\end{shortanswer}

\begin{shortanswer}{[GX-C01.S05-01]}
\(\dfrac{1}{4}\).
\end{shortanswer}

\begin{shortanswer}{[GX-C01.S05-02]}
\(5\sqrt{3}\).
\end{shortanswer}

\begin{shortanswer}{C01.S05 practica autonoma}
\begin{enumerate}[label=\textbf{PX-C01.S05-\arabic*:}, leftmargin=*]
  \item \(\dfrac{1}{9}\)
  \item \(3\sqrt{3}\)
  \item \(8\)
  \item \(3\)
  \item \(4\sqrt{5}\)
  \item \(20\)
\end{enumerate}
\end{shortanswer}

\begin{shortanswer}{[GX-C01.S06-01]}
\(\dfrac{5\sqrt{2}}{2}\).
\end{shortanswer}

\begin{shortanswer}{[GX-C01.S06-02]}
\(\sqrt{3}-\sqrt{2}\).
\end{shortanswer}

\begin{shortanswer}{C01.S06 practica autonoma}
\begin{enumerate}[label=\textbf{PX-C01.S06-\arabic*:}, leftmargin=*]
  \item \(\dfrac{2\sqrt{7}}{7}\)
  \item \(3+\sqrt{5}\)
  \item \(2-\sqrt{3}\)
  \item \(\sqrt{3}\)
  \item \(\dfrac{7\sqrt{7}+7}{6}\)
  \item \(-3-3\sqrt{2}\)
\end{enumerate}
\end{shortanswer}

\begin{shortanswer}{[GX-C01.S07-01]}
\(3\).
\end{shortanswer}

\begin{shortanswer}{[GX-C01.S07-02]}
\(125\).
\end{shortanswer}

\begin{shortanswer}{C01.S07 practica autonoma}
\begin{enumerate}[label=\textbf{PX-C01.S07-\arabic*:}, leftmargin=*]
  \item \(-2\)
  \item \(3\)
  \item \(5\)
  \item \(5\)
  \item \(1\)
  \item \(49\)
\end{enumerate}
\end{shortanswer}
